%------------------------------------------------------------------
%
% Vorlage für Abschlussarbeiten an der Technischen Hochschule Ingolstadt
% Bachelorarbeit/Masterarbeit
%
% Angepasst auf Seminararbeit
%
% V1 05.12.2011		Dr. Paul Spannaus
% V2 07.07.2012		Dr. Paul Spannaus
% V3 10.03.2013		Dr. Paul Spannaus
% V4 28.01.2014		Dr. Paul Spannaus
% V5 31.10.2014 	Dominik Schlecht
% 	
% -----------------------------------------------------------------
% -----------------------------------------------------------------

% Dokumentenklasse
\documentclass[11pt,DIV=11,twoside, openright, BCOR=3mm]{scrreprt} % 

\usepackage{amssymb}
\usepackage{array}
\usepackage{tabularx}
\usepackage[backend=biber,]{biblatex}
\usepackage[utf8]{luainputenc}

\renewcommand{\familydefault}{\sfdefault}
\usepackage{graphicx}

\usepackage{scrpage2} 	% Kopf & Fußzeile im KOMA Stil
\pagestyle{scrheadings}	% Aktiviert Verwendung vordefinierter Kolumnentitel
\clearscrheadfoot 		% alle Standard-Werte und Formatierungen löschen
\setkomafont{pagehead}{}	% Schriftart in Kopfzeile, \scshape = Kapitelchen
\automark[chapter]{section} % [linke Seite]{rechte Seite}
%\ohead{\def\pagestyle{PDTS}{\hrulefill\includegraphics[width = 6cm]{bilder/thi_logo_quer_cropped}}}
\ohead{\includegraphics[width = 3cm]{allianz_logo.jpg}}
\ihead{Fragebogen Penetrations-Test}

\setlength{\headsep}{7mm}		% Textabstand zur Kopfzeile
\setlength{\footskip}{7mm}		% Abstand zur Fußzeile

\ofoot{\vspace{-0.3cm} \pagemark} 
						
\setheadsepline{.2pt}
\setfootsepline{.4pt}	% Trennlinie Fußzeile und Textkörper

\newcolumntype{S}{>{\centering\arraybackslash}m{2em}}

\ifoot{\vspace{-1.7cm}
\begin{tabular}{l || l}
% Change me!
\textit{Unternehmen} 			& \textit{Pentester}\\ 
\textit{Straße Nr}			 	& \textit{Tel} \\
\textit{PLZ Ort} 					& \textit{E-Mail}\\
\textit{USTID}						& \textit{Webseite}
\end{tabular}
}



%%-------------Silbentrennung--------------------------------------
\hyphenation{}

%%-------------Index-----------------------------------------------
\makeindex
%\makenomenclature

%------------------------------------------------------------------
% Angabe der zu verwendenden Dateien, sodass nur z.B. das aktuelle 
% Dokument compiliert wird; den Rest mit % auskommentieren
\includeonly{%
	chapter/title_page,%
	chapter/thanks_statement,
	chapter/vorwort,
	chapter/appendix%
}


%------------------------------------------------------------------------

%-----------------------------------------------------------------
%---------------Dokumentenbeginn----------------------------------
%-----------------------------------------------------------------
	\begin{document}
			
			
% Titelseite
\begin{titlepage}

\phantom{tmpText}

\vspace{1cm}

\begin{figure}[h!]
\centering
\includegraphics[width=\textwidth]{bilder/thi_allianz_logo}
\end{figure}

  \begin{center}

%\vspace{1cm}
    $ $
    \\[3ex]
    \textbf{{\large Masterarbeit} \\[3ex]
    {\large Entwickung einer Umgebung für das Durchführen von}\\[1ex]
    {\LARGE Penetrationstests an mobilen Anwendungen} \\[1ex]
    {\large sowie die Klärung zugehöriger Prozesse}\\[6ex]
    zur Erlangung des akademisches Grades eines \\
    Master of Science\\[2ex]
    %
    \vfill
    %
    angefertigt von \\
    Dominik Gunther Florian Schlecht \\
    {\small \normalfont Matrikelnummer: 00032209}\\[2ex] %Vorname Nachname
    %
    \vfill
    %
    Betreuer} \\[2ex]
    \begin{tabular}{ll}
      Erstprüfer: & Prof. Hahndel \\
      Zweitprüfer: & Prof. von Koch \\
      Allianz Deutschland AG: & Herr Muncan und Herr Gerhager\\
      \\
      \\
      Fakultät: & Elektrotechnik und Informatik\\
      Studiengang: & Informatik\\
      Schwerpunkt: & Security \& Safety\\
      \\
      Abgabedatum: & 01. April 2016
    \end{tabular} \\[2ex]
    %
    \vfill
    %
    Ingolstadt, \today
  \end{center}
\end{titlepage}
			
				\cleardoublepage

			
%------------------------------------------------------------------------

% Nummerierung
\pagenumbering{roman}

% Sperrvermerk
%\chapter*{Sperrvermerk}
%Die vorliegende Arbeit \glqq Security Awareness Framework\grqq\ beinhaltet interne vertrauliche Informationen der Allianz Deutschland AG. Die Weitergabe des Inhalts der Arbeit im Gesamten oder in Teilen sowie das Anfertigen von Kopien oder Abschriften -- auch in digitaler Form -- sind grunds\"{a}tzlich untersagt. Ausnahmen bed\"{u}rfen der schriftlichen Genehmigung der Allianz Deutschland AG.
%\addcontentsline{toc}{chapter}{Sperrvermerk}

% Erkl\"{a}rung
\chapter*{Erkl\"{a}rung}
Hiermit erkl\"{a}re ich, dass ich die vorliegende Masterarbeit bis auf die offizielle Betreuung durch die Betreuer selbstst\"{a}ndig und ohne fremde Hilfe verfasst habe.\par
Die verwendeten Quellen sowie die verwendeten Hilfsmittel sind vollst\"{a}ndig angegeben. W\"{o}rtlich \"{u}bernommene Textteile und \"{u}bernommene Bilder und Zeichnungen sind in jedem Einzelfall kenntlich gemacht. \\[10ex]
Ingolstadt, \today
\addcontentsline{toc}{chapter}{Erkl\"{a}rung}

% Danksagung (optional)
%\chapter*{Danksagung}
%An dieser Stelle möchte ich mich bei allen Bedanken, die mich bei der Erstellung dieser Arbeit unterstützt haben. Besonderer Dank dabei gebührt den Kollegen aus der Allianz Deutschland AG und meinen Eltern.

%\begin{flushright}
%\sffamily Dominik Schlecht, \today
%\end{flushright}
%\addcontentsline{toc}{chapter}{Danksagung}


%------------------------------------------------------------------------

				\cleardoublepage	
					
			% Inhaltsverzeichnis
			%\setcounter{\secnumdepth}{2} % Zwei Ebenen im Inhaltsverzeichnis auflisten
			\setcounter{tocdepth}{4}
			\tableofcontents
				\cleardoublepage	
		
			% Abkrzungsverzeichnis
			%\printnomenclature
				%\cleardoublepage	
			
%------------------------------------------------------------------------
			% Neunummerierung des Hauptteils
			\pagenumbering{arabic}
			\setcounter{page}{1}            
            \cleardoublepage
            
            
            %\chapter{Einleitung}
Smartphones verbreiten sich immer stärker. Dies Zeigt auch eine Statistik von Gardner, nach welcher die Absätze von Mobilen Geräten die von herkömmlichen Rechnern und Laptops weiter übertreffen werden \cite{GartnerSales}, siehe Tabelle \ref{ref:GartnerSalesTable}. Mit diesem Fortschritt im Bereich der Verkäufe steigt natürlich auch die Nutzung von Smartphones. Diese können genutzt werden um im Internet zu "surfen", Medien-Inhalte wiederzugeben oder zu Chatten. Dies sind im Bezug auf Datenschutz und Risiko relativ ungefährliche Anwendungen. Jedoch können auch kritischere Handlungen vollzogen werden, wie zum Beispiel Online-Banking oder das Verwalten von Versicherungs-Verträgen. Das dies gemacht wird, ist eine notwendige Konsequenz aus der Natur des Smartphones und deren User. So ist es einfach praktisch, seinen Kontostand schnell von unterwegs einsehen zu können. Jedoch müssen solche Apps dann so gestaltet sein, dass ein Missbrauch nicht oder nur mit unverhältnismäßig hohem Aufwand möglich ist. Dies ist die Aufgabe der Unternehmen, welche solche Apps bereit stellen. Dies stellt viele vor unerwartet hohe Herausforderungen. Mussten bisher nur lokale oder Web-Anwendungen getestet werden, sind es nun Applikationen auf mobilen Geräten mit einem oft wesentlich höheren Anteil an Web-Services und APIs.

\begin{figure}[htbp]
	\centering
	\begin{tabular}{ l r r r r}
		Device Type & 2015 & 2016 & 2017 & 2018 \\ \hline
		Traditional PCs (Desk-Based and Notebook) & 244 & 228 & 223 & 216 \\
		Ultramobiles (Premium) & 45 & 57 & 73 & 90 \\
		PC Market & 289 & 284 & 296 & 306 \\
		Ultramobiles (Basic and Utility) & 195 & 188 & 188 & 194 \\
		Computing Devices Market & 484 & 473 & 485 & 500 \\
		Mobile Phones & 1,917 & 1,943 & 1,983 & 2,022 \\
		Total Devices Market & 2,401 & 2,416 & 2,468 & 2,521 \\
	\end{tabular}
	\label{ref:GartnerSalesTable}
	\caption{Worldwide Devices Shipments by Device Type, 2015-2018 (Millions of Units)\cite{GartnerSales}}
\end{figure}

Die Allianz Deutschland AG ist ein solches Unternehmen. Bisher waren viele Wege zum Kunden analog, also per Brief. Es gab nicht viele elektronische Schnittstellen. Vor 2 Jahren wurde eine Digitalisierungs-Strategie gegenüber beschlossen und die Web-Anwendung "Meine-Allianz" entwickelt. In dieser können Versicherungsnehmer nicht nur Verträge anschauen, sondern auch Änderungen an diesen Vornehmen. Sollte diese Anwendung eine schwerwiegende Schwachstelle aufweisen, hätten dies für die Allianz nicht nur Schäden in Bezug auf die Reputation, sondern eventuell zusätzliche rechtliche Folgen, da Krankendaten unter §203 StGB fallen. Um dies zu verhindern, sind Prozesse in der Anwendungsentwicklung notwendig, welche die Sicherheit einer Anwendung sicherstellen.

Diese Prozesse umfassen Komponenten wie Source-Scanning, Penetrations-Test sowie eine regelmäßig wiederkehrende Prüfung von Anwendungen, selbst nach dem Release.

In dieser Arbeit werden als Hinleitung allgemeine Prozesse und Fakten rund um die Applikations-Sicherheit erläutert. Im weiterführenden Teil wird ein bereits bestehendes Open-Source-Tool zum automatischen Testen mobiler Anwendungen um Features wie TODO erweitert.


            \chapter{Einleitung}
            \chapter{Arten von Penetrationstests}
	            \section{Web-Application}
	            \section{Wireless}
	            \section{Social-Engineering}
	            \section{Mobile-Applications}
	            \section{Blackbox/Whitebox}
            \chapter{Prozesse zu Penetrationstests}
	            \section{Vorbereitung}
		            \subsection{Aufwandsschätzung}
		            \subsection{Rechtliche Aspekte}
		            	\subsubsection{NDA}
		            	\subsubsection{Haftungsausschluss}
		            	\subsubsection{Absicherung gegen §203STGB}
		            \subsection{Technische Aspekte}
		            	\subsubsection{Infrastruktur}
		    			\subsubsection{Tools}
	            \section{Durchführung}
	            	\subsection{Bewertung von Findings}
	            		\subsubsection{CVSS}
	            		\subsubsection{Alternative Modelle}
					\subsection{Dokumentation}	            	
	            	\subsection{Reporting}    	
	            \section{Nachbereitung/Report}
	            	\subsection{Inhalte}
	            	\subsection{Vorlage}
	     	\chapter{Penetrationstests mobiler Anwendungen}
	     		\section{Aktuelle Situation und Vergleich}
	     			\subsection{IOS}
	     			\subsection{Windows-Phone}
	     			\subsection{Android}
	     		\section{Anforderungen}
	     		\section{Laboraufbau}
	     		\section{Entwicklung der Umgebung}
	     			\subsection{Technisches Detail 1}
	     			\subsection{Technisches Detail 2}
	     		\section{Abgleich mit Anforderungen}
	     	\chapter{Anwendung der Umgebung}
	     		\section{Pentest Anwendung 1}
	     		\section{Pentest Anwendung 2}
	     	\chapter{Fazit}
                \cleardoublepage	
	
%------------------------------------------------------------------------
			
			% Anhang
			
%------------------------------------------------------------------------
% Anhänge
\appendix
\chapter{Quellcode}

\section{Pentest-Helper}
\subsection{Latex}\label{ap:FragTex}
\lstset{language=TeX}
\lstinputlisting[caption={packages.tex}]{./Fragenkatalog/packages.tex}
\lstinputlisting[caption={question\_templates.tex}]{./Fragenkatalog/question_templates.tex}
\lstinputlisting[caption={Frage.tex}]{./Fragenkatalog/Fragen.tex}

\subsection{Docker}\label{ap:Docker}
\lstinputlisting[caption={Dockerfile für Entwicklungsversion}]{./logs/Dockerfile}

\section{iOS-Simulator}
\subsection{Control-Script}\label{ap:simcontrol}
\lstset{language=Python}
\lstinputlisting[caption={control.py}]{./logs/control.py}
\subsection{MITM-Server}\label{ap:mitmserver}
\lstset{language=Python}
\lstinputlisting[caption={MITM-Server}]{./logs/mitm_server.py}

\section{MobSF}
%\subsection{BinSkim}label{ap:BinSkimRegeln}
%\lstinputlisting[caption={BinSkim-Regeln}]{./logs/20170213_binskim_rules.txt}

\subsection{BinScope-Installer}\label{ap:BinScopeInstaller}
\lstset{language=Python}
\lstinputlisting[caption={BinScope-Installer}]{./logs/binscope_install.py}

%\newpage
%\subsection{SAF.cgi}
%\lstinputlisting[caption={SAF.cgi}]{../../SAF/SAFConsole/web/cgi-bin/SAF.cgi}

%\subsection{framework}\label{code:framework}
%\lstinputlisting[caption={SAFSubRout.pm}]{../../SAF/SAFConsole/framework/SAFSubRout.pm}


%\newpage
%\subsection{HTML-Snippets}
%\lstset{language=HTML}
%\lstinputlisting[caption={htmlStart.snip}]{../../SAF/SAFConsole/web/cgi-bin/snippets/htmlStart.snip}
%\lstinputlisting[caption={htmlEnd.snip}]{../../SAF/SAFConsole/web/cgi-bin/snippets/htmlEnd.snip}

%\section{Module}
%\subsection{Cracking}
%\subsubsection{John the Ripper}
%\lstset{language=Perl}
%\lstinputlisting[caption={john.mod}]{../../SAF/SAFConsole/modules/cracking/john.mod}
%\lstinputlisting[caption={johnDemo.pl}]{../../SAF/SAFConsole/modules/cracking/john/johnDemo.pl}

%\subsection{Sniffing}
%\subsubsection{Ettercap-Bridged}
%\lstset{language=Perl}
%\lstinputlisting[caption={ettercap\_bridged.mod}]{../../SAF/SAFConsole/modules/sniffing/ettercap_bridged.mod}
%\lstinputlisting[caption={ettercap\_bridged\_demo.pl}]{../../SAF/SAFConsole/modules/sniffing/ettercap_bridged/ettercap_bridged_demo.pl}

%TODO Module, Config
			
			% Literaturliste
			\newpage
			%\bibliographystyle{plaindin}	% DIN Norm für Literaturdarstellung
			%\bibliography{ref/ref_liste}	% Pfad und Datei der Ref-Datenbank
			\printbibliography % Literaturliste
		
	\end{document}
%-----------------------------------------------------------------
